\section{Dynamic recovery after an initially available route}
\label{sec:dynamic-recovery}
This separately frozen exploratory pilot distinguishes route availability at
feedback, authorization at execution, and the eventual business effect. It reuses
the bound Host and independent oracle without modification, and two preselected
checkpoint definitions (FORMAL-G-003 and FORMAL-G-004) from the earlier feedback
study. Each initially has a valid reviewed candidate, active grant and available
reuse, reauthorization and fresh-proposal routes. The initial visible state is
identical across the four conditions of a task. Eight hosted continuations, eight
adaptive-rule continuations and eight cached-route controls were executed on
23 October 2026 after freezing the protocol, inputs, code, allocation and limits.
They are not added to the original 32 G episodes.

Immediately after the controller chooses its first continuation commit, but before
the Host executes it, the harness either leaves state unchanged, revokes the grant
named by that action, advances the predecessor, or revokes active grants and
disables the reviewer. Selected-grant revocation is an adaptive, controlled
intervention, not a sample of naturally occurring administrative changes.
After the triggered return, all three controllers receive the same current
administrative facts. In the predecessor condition, the existing Host advances the
version while preserving the current reference-note value; a fresh candidate and
authorization can therefore restore a legitimate path. The unavailable-reviewer
condition has no legal continuation route.

All eight hosted first commits were admissible before intervention; only the two
static commits remained admissible at execution. The two revocations and two
unavailable-reviewer cases returned \texttt{AUTHORIZATION\_INVALID}; the two
predecessor changes returned \texttt{STALE\_VERSION}. These six deterministic
rejections are host safety behavior, not autonomous-recovery successes.

\begin{table}[htbp]
\centering\small
\caption{Dynamic pilot: independently read-back business recovery. Each controller
has two trajectories per condition; the controllers are not pooled.}
\begin{tabular}{lrrr}
\toprule
Condition & Hosted & Adaptive rule & Cached route\\
\midrule
Static & 2/2 & 2/2 & 2/2\\
Selected grant revoked & 2/2 & 2/2 & 0/2\\
Predecessor advanced & 2/2 & 2/2 & 0/2\\
No legal route & 0/2 & 0/2 & 0/2\\
\bottomrule
\end{tabular}
\end{table}

The hosted and adaptive-rule controllers reauthorized after revocation and rebuilt
the candidate and authorization after a predecessor advance. All three controllers
correctly stopped in both no-legal-route cases, with unchanged target values and
accurate observed reports. The cached-route controller also stopped after the
other dynamic rejections, leaving four feasible corrections incomplete. A fresh
SQLite connection independently read the persisted host state: all 24 trajectories
had integrity one, and all eight hosted final reports matched that state. Negative
controls rejected a forged success report and unrelated-field damage, and returned
unknown when independent readback was removed. All 24 saved trajectories were
reproduced exactly by offline replay.

\paragraph{A preserved measurement limitation.}
The frozen scorer additionally defines \texttt{task\_completion} by requiring a
successful \texttt{get\_state} at the final version, alongside business recovery
and an accurate report. This \emph{state-query-confirmed completion} measure is
narrower than the original task's allowance for receipt or state evidence. It is
0/6 in feasible hosted continuations: five instead queried a successful receipt,
and one used the successful commit response. All six business corrections are
independently verified and accurately reported. The frozen zeros are retained,
not interpreted as six failed business tasks. Adaptive rules score 6/6 and cached
routes 2/6 on this narrower measure. The pilot thus also exposes a measurement
dependence on confirmation-tool choice; subsequent studies should measure
completion from independent authoritative state and truthful reporting, with the
query choice reported separately.

The hosted collection used 35 requests and 27 continuation tool calls; adaptive
rules used 28 calls and cached routes 16. All responses returned the alias
\texttt{deepseek-v4-flash} and ended naturally, with no truncation, invalid output,
transport retry or budget exhaustion. Eight decisions, a 32,768-token output cap,
120 seconds per request and 900 seconds per episode were fixed before collection.
The raw provider-usage fields and local elapsed times are retained in
\texttt{Dynamic-recovery-reproduction.zip}; they do not estimate production latency.
Two reused task definitions, one realization per cell, controlled interventions
and a mutable service alias do not support statistical generalization, a
feedback-format comparison or LLM superiority. Matching adaptive-rule outcomes
show that the fixture is solvable without an LLM. The evidence establishes a
concrete temporal separation between an available route and a subsequently
authorized commit, and between gate rejection and actual recovery.
